% ============================================================================
% PROPOSED LaTeX INSERTS — threshold-vs-ranking decomposition
% Framing: grows out of the AUC/accuracy observation already stated in
% Section VI-F (subsec:feature-size); presented as a planned mechanistic
% analysis + an operational recipe, NOT as a rebuttal.
% ============================================================================


% ----------------------------------------------------------------------------
% INSERT 0 — CONSISTENCY FIX in Section VI-F (subsec:feature-size).
% REQUIRED if Insert 1 is added: otherwise VI-F frames threshold transfer as the
% sole cause of the S->M collapse, contradicting the ~55/45 split shown later.
%
% In the paragraph beginning "\textsc{auc} values follow the same overall pattern
% as accuracy.", REPLACE the sentence:
%   "While \textsc{auc} remains relatively high even at larger values of $K$,
%    accuracy in the \textsc{static}~$\rightarrow$~\textsc{mobile} direction
%    degrades substantially, indicating that decision thresholds learned in one
%    configuration do not transfer effectively to the other."
% WITH:

While \textsc{auc} remains relatively high even at larger values of $K$, accuracy
in the \textsc{static}~$\rightarrow$~\textsc{mobile} direction degrades
substantially, indicating that part of the gap reflects a decision threshold that
does not transfer across configurations. Section~\ref{subsec:threshold-decomposition}
quantifies this: threshold transfer accounts for roughly half of the degradation,
while a comparable component persists even under an oracle threshold, reflecting a
genuine loss of cross-domain separability.


% ----------------------------------------------------------------------------
% INSERT 0b — ABSTRACT (optional, softened; no numbers).
% Append after the existing sentence:
%   "Restricting the model to a compact invariant feature subset of four metrics
%    yields near-symmetric cross-domain transfer ($\approx 0.86$ in both
%    directions)."

Analysis indicates that the cross-domain degradation arises from both reduced class
separability and decision-threshold transfer, the latter being largely recoverable
through lightweight target-domain recalibration.


% ----------------------------------------------------------------------------
% INSERT 1 — new Results subsection.
% Place immediately AFTER \subsection{Feature Ablation and Augmentation ...}
% (subsec:feature-ablation-and-augmentation), i.e. as the last subsection of
% Section VI, just before \section{Discussion}.
% ----------------------------------------------------------------------------

\subsection{Threshold Transfer versus Feature Separability}
\label{subsec:threshold-decomposition}

The feature-size analysis (Section~\ref{subsec:feature-size}) observed that
cross-domain \textsc{auc} remains high in regimes where accuracy collapses: at
$K=33$ the \textsc{static}~$\rightarrow$~\textsc{mobile} \textsc{auc} is $0.8348$
against an accuracy of $0.6718$, whereas \textsc{Universal-4} attains $0.8974$ and
$0.8639$. Because \textsc{auc} is threshold-independent while accuracy is not, this
divergence indicates that part of the accuracy gap reflects a decision threshold
that does not transfer across configurations, rather than a loss of separability.
We make this quantitative by decomposing the gap into a \emph{ranking} component,
attributable to feature separability, and a \emph{threshold} component,
attributable to calibration transfer.

The decomposition is performed on identical model scores. For each direction and
each $K\in\{4,33\}$ we use the pipeline of Section~\ref{subsec:feature-size}
(\textsc{catboost}, group-aware split, source-validation threshold) and, over the
same $20$ bootstrap seeds, hold a group-aware $40\%$ target fold out for
evaluation. On this fold we vary \emph{only} the decision threshold across three
regimes: (i) the operational \emph{source} threshold tuned on source validation;
(ii) an \emph{oracle} threshold tuned on a disjoint target-calibration fold; and
(iii) \emph{few-shot} recalibration, in which the threshold is tuned on a small
number of labeled target runs. Writing $\mathrm{Acc}(K,\text{regime})$ for the
target-fold accuracy, we define, per seed,
\begin{align}
\Delta_{\text{total}}     &= \mathrm{Acc}(4,\text{src}) - \mathrm{Acc}(33,\text{src}),\\
\Delta_{\text{ranking}}   &= \mathrm{Acc}(4,\text{oracle}) - \mathrm{Acc}(33,\text{oracle}),\\
\Delta_{\text{threshold}} &= \Delta_{\text{total}} - \Delta_{\text{ranking}},
\end{align}
so that $\Delta_{\text{ranking}}$ isolates the separability advantage that survives
optimal thresholding and $\Delta_{\text{threshold}}$ captures the residual loss due
to threshold transfer. Confidence intervals are paired across seeds and corroborated
by a run-level cluster bootstrap.

Table~\ref{tab:threshold-decomposition} reports the result. In the
\textsc{static}~$\rightarrow$~\textsc{mobile} direction the $0.192$ gap separates
into $\Delta_{\text{ranking}}=0.106\,[0.097,0.115]$ and
$\Delta_{\text{threshold}}=0.086\,[0.078,0.094]$ (both $p<10^{-14}$). Roughly
$55\%$ of the gap is a genuine separability advantage of \textsc{Universal-4}---and,
notably, it is \emph{larger} than the $0.063$ \textsc{auc} gap, indicating that the
compact subset improves accuracy at the operating point beyond what the ranking
summary alone reflects. The remaining $\approx 45\%$ is a calibration effect:
re-tuning the threshold alone, without altering a single feature, raises the
$33$-feature detector from $0.672$ to $0.759$. The achievable (calibration-fold)
oracle equals the in-fold optimum to four decimals, so this is a genuine
operating-point limit rather than an estimation artifact. The
\textsc{mobile}~$\rightarrow$~\textsc{static} direction, which does not collapse,
shows a negligible threshold component ($0.003$), confirming that the calibration
effect is specific to the asymmetric \textsc{static}~$\rightarrow$~\textsc{mobile}
transfer.

The mechanism is a shift in the score distribution under mobility. For $K=33$ the
\textsc{static}~$\rightarrow$~\textsc{mobile} accuracy-optimal threshold moves by
$0.36$ between domains (logit-space Wasserstein distance $1.28$ between the
source-test and target-test score distributions), so a single source-tuned
threshold lands far from the target optimum even though the ranking is largely
preserved. For \textsc{Universal-4} the same shift is only $0.06$ (Wasserstein
$0.61$): its score distribution is domain-stable, so its threshold transfers
unchanged. The defense signal is therefore not lost under the full feature set;
it is mis-located.

This decomposition also yields a practical recalibration recipe and, at the same
time, sharpens the separability claim. Because the threshold component is
recoverable from target labels alone, tuning only the threshold on $25$ labeled
target runs raises $33$-feature \textsc{static}~$\rightarrow$~\textsc{mobile}
accuracy from $0.672$ to $0.746$, with $500$ runs reaching the $0.759$ oracle and
isotonic recalibration ($n\geq250$) attaining $0.779$. Crucially, even after the
threshold has been fully corrected the recalibrated $33$-feature detector plateaus
near its oracle and remains about $0.105$ below \textsc{Universal-4} ($0.864$)---a
residual that coincides with $\Delta_{\text{ranking}}$. Threshold recalibration
therefore cannot close the gap on its own: \textsc{Universal-4} is not merely
better calibrated but genuinely more separable under domain shift, and it reaches
this operating point with no target labels at all.

\begin{table}[t]
\centering
\footnotesize
\caption{Decomposition of the cross-domain accuracy gap between
\textsc{Universal-4} ($K{=}4$) and the full feature set ($K{=}33$) into a
feature-separability (ranking) component and a threshold-transfer (calibration)
component. All quantities use identical model scores over the same $20$ bootstrap
seeds; only the decision threshold differs across regimes. Operating-point
accuracies are mean~$\pm$~std; gap components are mean~[$95\%$~CI].}
\label{tab:threshold-decomposition}
\begin{tabular}{lcc}
\toprule
 & S$\rightarrow$M & M$\rightarrow$S \\
\midrule
\multicolumn{3}{l}{\textit{Operating-point accuracy, $K=33$}}\\
\quad Source threshold & $0.6716 \pm 0.0179$ & $0.8419 \pm 0.0042$ \\
\quad Oracle threshold & $0.7588 \pm 0.0208$ & $0.8453 \pm 0.0035$ \\
\midrule
\multicolumn{3}{l}{\textit{Operating-point accuracy, $K=4$ (\textsc{Universal-4})}}\\
\quad Source threshold & $0.8638 \pm 0.0046$ & $0.8584 \pm 0.0029$ \\
\quad Oracle threshold & $0.8648 \pm 0.0040$ & $0.8590 \pm 0.0030$ \\
\midrule
\multicolumn{3}{l}{\textit{Gap decomposition ($K{=}4$ minus $K{=}33$)}}\\
\quad $\Delta_{\text{total}}$ (source thr.)   & $0.192\,[0.184,0.201]$ & $0.016\,[0.014,0.018]$ \\
\quad $\Delta_{\text{ranking}}$ (oracle thr.) & $0.106\,[0.097,0.115]$ & $0.014\,[0.012,0.015]$ \\
\quad $\Delta_{\text{threshold}}$ (residual)  & $0.086\,[0.078,0.094]$ & $0.003\,[0.001,0.004]$ \\
\bottomrule
\end{tabular}
\end{table}


% ----------------------------------------------------------------------------
% INSERT 2 — revise the over-attributing sentence in
% \subsection{Feature Invariance as the Determinant of Generalization}
% (subsec:invariance).
%
% (Optional) retitle to:
%   \subsection{Feature Invariance and Score-Distribution Stability as
%   Determinants of Generalization}
% ----------------------------------------------------------------------------

% REPLACE the final sentence of the first paragraph of subsec:invariance
% ("...and exceeds the full feature set by roughly $0.19$ absolute accuracy on
%  S$\rightarrow$M.") WITH:

and exceeds the full feature set by $0.192$ absolute accuracy on
S$\rightarrow$M. As shown in Section~\ref{subsec:threshold-decomposition}, this
$0.192$ gap is not a single effect: about $0.106$ is a genuine separability
advantage that persists under an oracle decision threshold (and exceeds the
$0.063$ \textsc{auc} gap), while about $0.086$ is a calibration effect---the
source-tuned threshold fails to transfer because the full-feature score
distribution drifts under mobility, whereas the \textsc{Universal-4} distribution
does not.

% REPLACE the closing principle paragraph of subsec:invariance
% ("The principle that emerges is straightforward. Cross-domain robustness ...
%  Features with weaker but more stable univariate separability generalize.")
% WITH:

The principle that emerges is straightforward, and it operates through two
complementary channels. Cross-domain robustness in this setting is governed by
feature invariance rather than by univariate discriminative strength in the source
domain: strong but configuration-dependent features, most prominently the delay-
and jitter-based metrics, fail to transfer, while features with weaker but more
stable univariate separability generalize. Invariance acts first by preserving the
\emph{ranking} of samples across domains, and second---previously implicit in the
high cross-domain \textsc{auc}---by preserving the \emph{location} of the score
distribution, so that a single decision threshold remains valid across
configurations. The full feature set satisfies the first property only partially
and the second hardly at all (its \textsc{static}~$\rightarrow$~\textsc{mobile}
optimal threshold shifts by $0.36$, versus $0.06$ for \textsc{Universal-4}); the
compact invariant subset satisfies both. Selecting invariant features therefore
yields not only separable but also \emph{calibrated} cross-domain decisions, which
is what allows \textsc{Universal-4} to transfer without any target-domain
recalibration.


% ----------------------------------------------------------------------------
% INSERT 3 (optional, subtle) — Conclusion.
% Append to the paragraph that introduces Universal-4's near-symmetric transfer
% ("...near-symmetric transfer in both directions ... at a minor cost in
%  source-domain accuracy.")
% ----------------------------------------------------------------------------

A decomposition of this gap on identical scores
(Section~\ref{subsec:threshold-decomposition}) attributes roughly $55\%$ of the
$K=33$ \textsc{static}~$\rightarrow$~\textsc{mobile} shortfall to feature
separability and the remainder to a decision threshold that does not transfer
across domains. The calibration component is recoverable---as few as $25$ labeled
target runs recover most of it---whereas the separability component is not: even
under an oracle threshold the full feature set remains well below
\textsc{Universal-4}. The compact invariant subset avoids both effects
simultaneously, transferring without any target-domain recalibration because it is
both more separable and better calibrated across configurations.

% ============================================================================
% Optional figures (already generated):
%   results/threshold_decomposition/fig_score_drift.{png,pdf}
%       -> 2x2 score-distribution histograms with t*_src / t*_tgt marked.
%   results/threshold_decomposition/fig_fewshot_recalibration.{png,pdf}
%       -> accuracy vs. number of labeled target runs.
% Either can be added to Section VI-H if a visual is desired; the text above is
% self-contained without them.
% ============================================================================
